\documentclass[../DoAn.tex]{subfiles}
\begin{document}

\begin{center}
    \Large{\textbf{TÓM TẮT NỘI DUNG ĐỒ ÁN}}\\
\end{center}
\vspace{1cm}

Bài toán bao phủ bản đồ là một nhiệm vụ quan trọng trong các hệ thống robot tự hành, đặc biệt khi robot cần đi qua toàn bộ một khu vực để phục vụ các nhiệm vụ có yêu cầu an toàn cao. Trong đồ án này, định hướng ứng dụng chính được đặt vào bài toán robot dò mìn và rà soát vật nguy hiểm trong một khu vực có rủi ro. Với dạng nhiệm vụ này, robot cần bao phủ càng đầy đủ khu vực làm việc càng tốt, đồng thời phải xét đến giới hạn năng lượng, địa hình di chuyển không đồng nhất, vật cản tĩnh, vật cản động và khả năng quay về căn cứ để sạc pin hoặc kết thúc nhiệm vụ an toàn. Trong đồ án này, căn cứ được hiểu là ô xuất phát của robot và cũng là nơi robot quay về để sạc pin. Ngoài bài toán dò mìn, mô hình cũng có thể được điều chỉnh để phục vụ các nhiệm vụ bao phủ khác như lau nhà tự động, kiểm tra kho bãi, cứu trợ cứu nạn hoặc khảo sát môi trường sau thảm họa.

Đồ án xây dựng một hệ thống mô phỏng robot bao phủ bản đồ lưới có xét năng lượng và vật cản động. Môi trường được biểu diễn bằng bản đồ lưới có trọng số, trong đó mỗi ô hợp lệ có thể mang chi phí di chuyển khác nhau nhằm mô phỏng các mức độ khó dễ của địa hình. Quá trình lập kế hoạch đường đi được xây dựng trên Dijkstra có hướng, với trạng thái bao gồm cả vị trí và hướng quay của robot, cho phép mô hình hóa riêng năng lượng di chuyển và năng lượng quay. Trước khi tiếp tục bao phủ một ô mục tiêu, robot kiểm tra khả năng quay về căn cứ, nhờ đó hạn chế trường hợp đi sâu vào khu vực làm việc nhưng không đủ năng lượng để trở về. Vật cản động được mô phỏng như các tác nhân chuyển động hợp tác; chúng không truy đuổi hay cố tình chặn robot, nhưng có thể làm thay đổi trạng thái an toàn của các ô trên bản đồ theo thời gian.

Kết quả của đồ án là một chương trình mô phỏng bằng C++ và OpenCV, có khả năng trực quan hóa quá trình di chuyển, ghi log hành vi, xuất dữ liệu benchmark và hỗ trợ so sánh giữa mô hình robot có chi phí quay với Metric H, trong đó chi phí quay được xem là không đáng kể. Các kết quả thực nghiệm cho thấy mô hình năng lượng quay không chỉ làm thay đổi tổng năng lượng tiêu thụ, mà còn ảnh hưởng trực tiếp tới hành vi bao phủ, số bước đi lại, số lần quay về sạc và hiệu suất bao phủ trong từng chu kỳ làm việc.

\begin{flushright}
Sinh viên thực hiện\\
\begin{tabular}{@{}c@{}}
\textit{Hà Kim Minh}
\end{tabular}
\end{flushright}

\end{document}
